\section{Discussion}
\label{sec:discussion}
% This section discusses the results of the data collected during the tests and the threats to the validity of this study.

% After data collection, it was possible to observe significant differences across the analyzed metrics. The results show that both rendering strategies performed well with smaller datasets. However, as the data size increased, application performance degraded. SSR achieved better results in the TTI and LCP metrics, while CSR yielded better outcomes in SI and FCP. Regarding JS Heap memory usage, there were no statistically significant differences, as shown in Table~\ref{tabela:estatistica}.

While both rendering strategies demonstrated good performance with smaller datasets, significant differences emerged as data volume increased, directly impacting application efficiency. Specifically, SSR consistently yielded better results for LCP and TTI , likely due to the server delivering a fully rendered HTML page, allowing the browser to display core content and become interactive more rapidly without extensive client-side JavaScript processing. Conversely, CSR exhibited superior performance in FCP and SI, especially with larger datasets where its ability to dynamically render content once initial scripts are loaded minimized the time to render visible elements. This distinction suggests that for content-heavy applications where initial load and SEO are paramount, SSR may be more advantageous, whereas highly interactive applications might benefit from CSR's dynamic rendering capabilities after initial asset loading. Regarding \textit{JS Heap} memory usage, there were no statistically significant differences, as shown in Table~\ref{tabela:estatistica}.

This study supports the findings of \citet{meredova2023ssr}, which analyzed server-side rendering (SSR) capabilities using the React and Vue front-end libraries. According to the study, the SSR approach enhances the user experience by reducing the initial page load time. This advantage makes SSR suitable for devices with slower internet connections and for applications where fast content loading is essential.
